\section{The Standard Instrumental Variable Estimator}
\label{sec:standard_iv}

\begin{definition}[Standard IV estimator]
    Consider the setup given in \ref{eq:structural}. The standard IV estimator $\hat{\beta}_{\mathrm{IV}}$ is given by
\[
    \hat{\beta}_{\mathrm{IV}} = \frac{\bar{S}_{ZY}}{\bar{S}_{ZX}}, \qquad \bar{S}_{ZY} = \frac{1}{n}\sum_{i=1}^{n}Z_iY_i, \quad \bar{S}_{ZX} = \frac{1}{n}\sum_{i=1}^{n}Z_iX_i.
\]
\end{definition}

\begin{lemma}[Error decomposition]
\label{lem:iv_decomp}
Under \eqref{eq:structural} and \ref{ass:exog},
\begin{equation}
\label{eq:iv_error}
    \hat{\beta}_{\mathrm{IV}} - \beta = \frac{\bar{S}_{Z\varepsilon}}{\bar{S}_{ZX}},
\end{equation}
where $\bar{S}_{Z\varepsilon} = \frac{1}{n}\sum_{i=1}^{n}Z_i\varepsilon_i$ satisfies $\E[\bar{S}_{Z\varepsilon}] = 0$.
\end{lemma}

\begin{proof}
Multiply both sides of $Y_i = \beta X_i + \varepsilon_i$ by $Z_i$ and average:
\[
    \bar{S}_{ZY} = \frac{1}{n}\sum_{i=1}^{n}Z_i(\beta X_i + \varepsilon_i) = \beta\bar{S}_{ZX} + \bar{S}_{Z\varepsilon}.
\]
Dividing by $\bar{S}_{ZX}$ gives $\hat{\beta}_{\mathrm{IV}} = \beta + \bar{S}_{Z\varepsilon}/\bar{S}_{ZX}$. By \ref{ass:exog}, $\E[Z_i\varepsilon_i] = 0$, so $\E[\bar{S}_{Z\varepsilon}] = 0$ by linearity of expectation. Thus, $\hat{\beta}_{\mathrm{IV}}$ is unbiased.
\end{proof}

\begin{theorem}[Standard IV concentration]
\label{thm:iv}
Assume \ref{ass:exog}--\ref{ass:moments}. For any $\delta \in (0,1)$, if
\begin{equation}
\label{eq:iv_strength}
    n \geq \frac{8\sigma_{ZX}^2}{\delta\,\mu_{ZX}^2},
\end{equation}
then
\begin{equation}
\label{eq:iv_bound}
    \Prob\!\left[\abs{\hat{\beta}_{\mathrm{IV}} - \beta} > \frac{2\sqrt{2}\,\sigma_{Z\varepsilon}}{\abs{\mu_{ZX}}\sqrt{\delta\, n}}\right] \leq \delta.
\end{equation}
\end{theorem}


\begin{proof}
The proof uses the two-event strategy applied to the full-sample decomposition~\eqref{eq:iv_error}.

\medskip\noindent\textbf{Step 1 (Two-event decomposition).} For a threshold $t > 0$ to be determined, define
\[
    A = \left\{\abs{\bar{S}_{ZX} - \mu_{ZX}} \leq \frac{\abs{\mu_{ZX}}}{2}\right\}, \qquad B = \left\{\abs{\bar{S}_{Z\varepsilon}} \leq \frac{t\abs{\mu_{ZX}}}{2}\right\}.
\]
\todo{Why?}
On event $A$, the reverse triangle inequality gives $\abs{\bar{S}_{ZX}} \geq \abs{\mu_{ZX}} - \abs{\bar{S}_{ZX} - \mu_{ZX}} \geq \abs{\mu_{ZX}}/2$. On $A \cap B$,
\[
    \abs{\hat{\beta}_{\mathrm{IV}} - \beta} = \frac{\abs{\bar{S}_{Z\varepsilon}}}{\abs{\bar{S}_{ZX}}} \leq \frac{t\abs{\mu_{ZX}}/2}{\abs{\mu_{ZX}}/2} = t.
\]

\medskip\noindent\textbf{Step 2 (Chebyshev on each event).} Since $\Var(\bar{S}_{ZX}) = \sigma_{ZX}^2/n$ and $\Var(\bar{S}_{Z\varepsilon}) = \sigma_{Z\varepsilon}^2/n$, Chebyshev's inequality gives
\[
    \Prob[A^c] \leq \frac{\sigma_{ZX}^2/n}{\mu_{ZX}^2/4} = \frac{4\sigma_{ZX}^2}{n\mu_{ZX}^2}, \qquad
    \Prob[B^c] \leq \frac{\sigma_{Z\varepsilon}^2/n}{t^2\mu_{ZX}^2/4} = \frac{4\sigma_{Z\varepsilon}^2}{nt^2\mu_{ZX}^2}. 
\]

\medskip\noindent\textbf{Step 3 (Union bound and solve for $t$).} We allocate $\delta/2$ to each event. The condition~\eqref{eq:iv_strength} ensures $\Prob[A^c] \leq \delta/2$. Setting $\Prob[B^c] \leq \delta/2$ and solving for $t$:
\[
    \frac{4\sigma_{Z\varepsilon}^2}{nt^2\mu_{ZX}^2} = \frac{\delta}{2} \implies t = \frac{2\sqrt{2}\,\sigma_{Z\varepsilon}}{\abs{\mu_{ZX}}\sqrt{\delta n}}.
\]
By the union bound, $\Prob\left[\abs{\hat{\beta}_{\mathrm{IV}} - \beta} > t\right] \leq \Prob[A^c] + \Prob[B^c] \leq \delta/2 + \delta/2 = \delta$.
\end{proof}

\begin{remark}[Polynomial dependence on $\delta$]
\label{rem:iv_polynomial}
The bound~\eqref{eq:iv_bound} depends on $\delta$ through the factor $1/\sqrt{\delta}$, which is polynomial. \cite{devroye2016sub} show that for any estimator based solely on the empirical mean of i.i.d.\ observations with finite variance, polynomial dependence on $\delta$ is unavoidable without additional distributional assumptions.
\end{remark}

\begin{remark}[Instrument strength depends on $\delta$]
The condition \eqref{eq:iv_strength} involves $\delta$: achieving higher confidence (smaller $\delta$) requires more samples to control the denominator. This dependence is polynomial in $1/\delta$.
\end{remark}

